% !TEX root = ../main.tex
% !TEX program = lualatex
\documentclass[../main.tex]{subfiles}
\IfSubfilesClassLoaded{\externaldocument{\subfix{../build/main}}}{}
% =============================================================================
% APPENDIX: COMBATANT COMMANDS
% DESCRIPTION: Geographic and functional combatant commands (CCMDs) reference.
% =============================================================================
\begin{document}
\chapter{Combatant Commands Reference}\label{app:combatant-commands}

\section{Command List}
\subsection{Summary}
\begin{description}[style=nextline, labelsep=0.5em, font=\bfseries]
  \item[Purpose.] This appendix is a board-speed list of the unified \acp{ccmd} and the geographic/functional split. Use it for recall, then return to the command-chain chapter for \ac{ccdr} authority and assigned-force discussion.
  \item[Current count.] The Department's public combatant-command page identifies 11 \acp{ccmd}, each with a geographic or functional mission for command and control in peace and war~\autocite{DOW-CombatantCommands-2026}.
  \item[Geographic vs.\ functional.] The public Department resource page describes geographic commands as having clearly delineated operating areas and functional commands as operating worldwide across geographic boundaries to provide unique capabilities~\autocite{DOW-UnifiedCombatantCommands-2026}.
  \item[Space Command cue.] Treat \ac{usspacecom} as the geographic/astrographic command: public Department reporting describes it as a geographic combatant command with an area of responsibility beginning 100 km above mean sea level and extending higher~\autocite{DoD-Spacecom-Strategic-Environment-2019}.
\end{description}

\noindent\textbf{Area of responsibility:} Do not use the guide's \texttt{aor} acronym key here; in this guide, \ac{aor} means Accumulated Operating Result for \ac{nwcf}. In combatant-command context, area of responsibility means the geographic or astrographic area assigned to a \ac{ccdr} for missions and forces~\autocite{USC-162,DOW-UnifiedCombatantCommands-2026}.

\textbf{Geographic (7)~\autocite{DOW-CombatantCommands-2026,DOW-UnifiedCombatantCommands-2026,DoD-Spacecom-Strategic-Environment-2019}:}
\begin{itemize}
  \item \ac{usafricom}
  \item \ac{uscentcom}
  \item \ac{useucom}
  \item \ac{usindopacom}
  \item \ac{usnorthcom}
  \item \ac{ussouthcom}
  \item \ac{usspacecom}\footnote{\ac{usspacecom} is treated as a geographic/astrographic command; public Department reporting describes its area of responsibility as beginning 100 km above mean sea level and extending higher~\autocite{DoD-Spacecom-Strategic-Environment-2019}.}
\end{itemize}
\textbf{Functional (4)~\autocite{DOW-CombatantCommands-2026,DOW-UnifiedCombatantCommands-2026}:}
\begin{itemize}
  \item \ac{uscybercom}
  \item \ac{ussocom}
  \item \ac{usstratcom}
  \item \ac{ustranscom}
\end{itemize}

\subsection{Quick Review}
\begin{itemize}
  \item Count: 11 unified \acp{ccmd}; memorize 7 geographic and 4 functional.
  \item Geographic: USAFRICOM, USCENTCOM, USEUCOM, USINDOPACOM, USNORTHCOM, USSOUTHCOM, and USSPACECOM.
  \item Functional: USCYBERCOM, USSOCOM, USSTRATCOM, and USTRANSCOM.
  \item Board trap: \ac{usspacecom} belongs in the geographic/astrographic count, not the functional count.
  \item Acronym trap: spell out combatant-command area of responsibility in prose because this guide's \ac{aor} acronym key is reserved for \ac{nwcf} Accumulated Operating Result.
\end{itemize}

%====================
% End of file
%====================
\IfSubfilesClassLoaded{
  \RenewDocumentCommand{\entryname}{}{\textbf{\color{Modern} Acronym}}
  \RenewDocumentCommand{\descriptionname}{}{\textbf{\color{Modern} Definition}}
  \printnoidxglossary[
    type=\acronymtype,
    title=Acronyms,
    style=long-booktabs
  ]}{}
\end{document}
